\section{Environnement technique}

% ----------------------------------------------------------------
% 📌 À quoi sert cette section ?
% ----------------------------------------------------------------
% Cette partie décrit l’ensemble des outils, technologies, langages et environnements
% utilisés pour mener à bien le projet.
% Elle permet de cadrer le développement et d’assurer la reproductibilité.

% ----------------------------------------------------------------
% 🧭 Comment la remplir ?
% ----------------------------------------------------------------
% ➤ Listez les outils logiciels (IDE, bibliothèques, frameworks, etc.)
% ➤ Indiquez les langages de programmation, OS, navigateurs ou environnements serveur ciblés.
% ➤ Mentionnez les éventuelles contraintes matérielles ou logicielles.
% ➤ Vous pouvez structurer les infos par catégories (langages, frameworks, outils, etc.)

% ----------------------------------------------------------------
% 🧾 Exemple rédigé :
% ----------------------------------------------------------------
% Le projet sera développé à l’aide des technologies suivantes :
%
% \textbf{Langages et Frameworks} :
% \begin{itemize}
%   \item JavaScript / React pour l’interface web
%   \item HTML / CSS pour la structuration visuelle
%   \item SVG natif manipulé en DOM
% \end{itemize}
%
% \textbf{Outils de développement} :
% \begin{itemize}
%   \item VSCode (éditeur de code)
%   \item GitHub pour le versionning
%   \item React CodeMirror pour l’édition sécurisée du SVG
% \end{itemize}
%
% \textbf{Sécurité / Sanitation} :
% \begin{itemize}
%   \item DOMPurify pour filtrer le code SVG injecté
% \end{itemize}
%
% \textbf{Déploiement} :
% \begin{itemize}
%   \item Hébergement via Vercel
% \end{itemize}

% ----------------------------------------------------------------
% 📊 Exemple de structure par tableau
% ----------------------------------------------------------------
% \begin{table}[H]
% \centering
% \begin{tabular}{|l|l|}
% \hline
% Élément & Détail \\
% \hline
% Langage principal & JavaScript \\
% Framework front-end & React \\
% Librairie de code sécurisé & DOMPurify \\
% Environnement de dev & VSCode + Git \\
% Plateforme de déploiement & Vercel \\
% \hline
% \end{tabular}
% \caption{Résumé de l’environnement technique}
% \end{table}

% ----------------------------------------------------------------
% ✍️ À vous de jouer : complétez ci-dessous
% ----------------------------------------------------------------

% (Suggestion de début)
L’environnement technique du projet repose sur les éléments suivants :

\begin{itemize}
  \item [À compléter...]
\end{itemize}
